% Vietnamese translation for OI Public Library
% Source-File: IOI中国国家候选队论文/2005/金恺/金恺.ppt
\documentclass[11pt]{article}
\usepackage{oipl-vn}

\title{Bản trình chiếu: Ghi chú hình học của Jin Kai}
\author{Jin Kai}
\date{2005}

\begin{document}
\maketitle

\origtitle{Slide companion for Jin Kai's geometry presentation}
\sourcepath{IOI China candidate team papers / 2005 / Jin Kai / Jin Kai.ppt}

\section{Phạm vi bản dịch}

Tệp gốc chỉ có PowerPoint. Phần chữ đọc được chứa các biến \(n,m,x,y\), điểm
\(O(0,0)\), \(A(a,0)\), \(B(b,0)\), các hướng \texttt{ESWN}, giới hạn
\(N\le 1000\), và một số bộ dữ liệu ví dụ như \(5,5,1,1\).

\section{Mô hình nhận diện được}

Các ký hiệu cho thấy một bài toán hình học tọa độ hoặc di chuyển trên mặt
phẳng. Điểm \(O,A,B\) cùng các giá trị \(a,b\) có thể tạo hệ trục hoặc các
mốc biên; chuỗi hướng \texttt{ESWN} gợi chuyển động theo đông, nam, tây, bắc.

\section{Hướng xử lý}

Với ràng buộc \(N\le 1000\), nhiều cách hình học rời rạc có thể dùng được:
mô phỏng sự kiện, chia đoạn, hoặc xét giao của đường đi với các mốc tọa độ.
Các giá trị lớn như \(10000\) cho thấy vẫn cần tránh lưới đầy đủ nếu miền tọa
độ rộng.

\section{Ghi chú}

Bản dịch này không suy diễn một đề bài hoàn chỉnh khi văn bản slide không đủ.
Nó giữ lại các ký hiệu chắc chắn đọc được và mô tả hướng thuật toán hợp lý
cho nhóm công thức đó.

\end{document}
